\documentclass[11pt]{amsart}

\usepackage[T1]{fontenc}
\usepackage{lmodern}
\usepackage{microtype}
\usepackage{mathtools,amssymb,amsthm}
\usepackage[hidelinks]{hyperref}

\hypersetup{
  pdftitle={Pure powers 3^7 and 3^8 in the integer group determinant image of D_54}
}

\newtheorem{theorem}{Theorem}
\newtheorem{proposition}[theorem]{Proposition}
\newtheorem{lemma}[theorem]{Lemma}
\theoremstyle{remark}
\newtheorem*{remark}{Remark}

\DeclareMathOperator{\Res}{Res}
\newcommand{\ZZ}{\mathbb Z}
\newcommand{\Sset}{\mathcal S}

\title[Pure powers in the determinant image of $D_{54}$]
{Pure powers $3^7$ and $3^8$ in the integer group determinant image of $D_{54}$}

\date{}

\subjclass[2020]{11C20, 11R06, 15B36, 20C05}
\keywords{integer group determinant, Lind--Lehmer problem, dihedral group,
cyclotomic resultant, group ring}


\begin{document}

\begin{abstract}
Let $\Sset(G)$ be the set of nonzero integer group determinants of a finite
group $G$. We exhibit two elements
$(x+x^4)+(x+x^4-1)s$ and
$(1+x+x^2+x^4+x^5)+(x+x^2+x^4+x^5)s$
of $\ZZ[D_{54}]$ whose group determinants are exactly $3^7=2187$ and
$3^8=6561$, with no cofactor: pure powers of $3$ with $v_2=0$. This is the new
content of the note. Granting the necessity and achievability statements of
Theorem 5.4 of Boerkoel and Pinner \cite{BP}, which are quoted here and not
reproved, these two pure powers close the exponent range $7\le b\le8$ left
between those statements at $p=3$, $k=3$ and yield the description
$\Sset(D_{54})=\{2^a3^bm:\gcd(m,6)=1,\ a=0\text{ or }a\ge2,\ b=0\text{ or }
b\ge7\}$, $m$ ranging over all of $\ZZ$.
\end{abstract}

\maketitle

\section{The window in Theorem 5.4}

Let $G$ be a finite group of order $N$ and let
$D_G(a)=\det\bigl(a_{uv^{-1}}\bigr)_{u,v\in G}$ be the group determinant of
$a=\sum_{g}a_g\,g\in\ZZ[G]$, an $N\times N$ integer determinant. Write
\[
 \Sset(G)=\{D_G(a):a\in\ZZ[G]\}\setminus\{0\}
\]
for the set of nonzero \emph{measures} of $G$. Two conventions are inherited
from \cite{BP}: the value $0=D_G(0)$, attained for every $G$, is excluded
throughout, and cofactors written $m$ range over \emph{all} of $\ZZ$ subject to
the stated coprimality, not merely the positive integers.
For the dihedral group
\[
 D_{2n}=\langle r,s \mid r^n=s^2=1,\ srs=r^{-1}\rangle
\]
an element of $\ZZ[D_{2n}]$ is written $f(r)+g(r)s$ with
$f,g\in\ZZ[x]/(x^n-1)$.

Boerkoel and Pinner \cite{BP} proved the following.

\begin{quote}
\textbf{Theorem 5.4.} For $G=D_{2p^2}$ with $p=3,5$ or $7$ the measures take
the form $2^ap^bm$, $(m,2p)=1$ with $a=0$ or $a\ge2$, and $b=0$ or $b\ge5$.

In general, if $G=D_{2p^k}$ with $p$ an odd prime and $k\ge2$ the measures must
take the form $2^ap^bm$, $(m,2p)=1$, with $a=0$ or $a\ge2$, and $b=0$ or
$b\ge2k+1$. We can achieve everything of this form with $b=0$ or $b\ge3k$.
There are measures with $b=2k+1$.
\end{quote}

Necessity is Lemma 4.4 of \cite{BP} (p.~9): if $p^\alpha\,\|\,n$ and $p\mid D$
then $p^{2\alpha+1}\mid D$. Achievability is assembled in the proof of Theorem
5.1 of \cite{BP} (pp.~15--16): Theorem 4.3 there gives every $m$ coprime to
$2p$, two displayed products give every $2^a$ with $a\ge2$, and a family indexed
by $\ell\ge k$ gives the pure powers $p^{\ell+2k}$, that is every $p^b$ with
$b\ge3k$. That proof then remarks \emph{``For $k=1$ these coincide and we have a
complete description of the measures''}, so the window
\[
 2k+1\;\le\;b\;\le\;3k-1
\]
is not determined either way by Theorem 5.4 for every $k\ge2$ (it lies strictly
between the necessity bound and the achievability bound stated there). At $k=2$ it is the single
value $b=5$, and the first clause of Theorem 5.4 covers it for $p=3,5,7$. For
$p=3$ and $k=3$, that is $G=D_{54}$ with $n=27$, the window is $b\in\{7,8\}$.
Neither \cite{BP} nor its e-print poses a question or conjecture about this
window; no priority or novelty claim is made here, and no literature search
beyond \cite{BP} was carried out. What we prove is that $3^7$ and $3^8$ are
measures of $D_{54}$ with cofactor $1$; the last sentence of Theorem 5.4
already asserts that some measure with $b=2k+1=7$ exists, and it is pure powers
that the multiplicative assembly of \cite{BP} consumes.



\section{The two elements}

Throughout $n=27$, $\zeta_d$ is a primitive $d$-th root of unity,
$f^*(x)=f(x^{-1})$ in $\ZZ[x]/(x^{27}-1)$, and
\[
 M(h)=\prod_{j=0}^{26}h(\zeta_{27}^{\,j})=\Res(x^{27}-1,\,h)
\]
is the circulant determinant of $h$. For $n$ odd the reduction of \cite{BP} is
\begin{equation}\label{eq:red}
 D_{D_{2n}}\bigl(f+gs\bigr)=M\bigl(f f^*-g g^*\bigr),
 \qquad h:=f f^*-g g^*.
\end{equation}

\begin{theorem}\label{thm:main}
Let $G=D_{54}$.
\begin{enumerate}
\item[\textup{(W1)}] For $f=x+x^4$ and $g=f-1=-1+x+x^4$,
\[
 h=-1+x+x^4+x^{23}+x^{26},\qquad h(1)=3,
\]
\[
 \Res(\Phi_3,h)=\Res(\Phi_9,h)=\Res(\Phi_{27},h)=9,
\]
\[
 D_G(f+gs)=3\cdot9^3=2187=3^7 .
\]
\item[\textup{(W2)}] For $f=1+x+x^2+x^4+x^5$ and $g=f-1=x+x^2+x^4+x^5$,
\[
 h=1+(x+x^{26})+(x^2+x^{25})+(x^4+x^{23})+(x^5+x^{22}),\qquad h(1)=9,
\]
\[
 \Res(\Phi_3,h)=\Res(\Phi_9,h)=\Res(\Phi_{27},h)=9,
\]
\[
 D_G(f+gs)=9\cdot9^3=6561=3^8 .
\]
\end{enumerate}
In both cases the value is a \emph{pure} power of $3$: $v_2=0$ and the cofactor
after removing all factors $2$ and $3$ is $1$.
\end{theorem}

\section{The valuation argument}\label{sec:valuation}

Both witnesses have the shape $g=f-1$, and then
$g g^*=(f-1)(f^*-1)=f f^*-f-f^*+1$, so
\begin{equation}\label{eq:h}
 h=f+f^*-1,
 \qquad\text{i.e.}\qquad
 h_0=2a_0-1,\quad h_i=a_i+a_{27-i}\ (1\le i\le26).
\end{equation}
In particular $h$ is symmetric, $h_i=h_{-i}$, and $h(1)=2f(1)-1$ is odd.

\begin{lemma}\label{lem:sq}
Let $h\in\ZZ[x]/(x^{27}-1)$ be symmetric. Then for $m=1,2,3$ there is
$P_m\in\ZZ$ with $\Res(\Phi_{3^m},h)=P_m^2$, and
\begin{equation}\label{eq:val}
 v_3\bigl(D_{D_{54}}(f+gs)\bigr)
 = v_3\bigl(h(1)\bigr)+2\sum_{m=1}^{3}v_3(P_m).
\end{equation}
Moreover $3\mid h(1)$ if and only if $3\mid P_m$, for each $m$.
\end{lemma}

\begin{proof}
Symmetry gives $h(\zeta^{-1})=h(\zeta)=\overline{h(\zeta)}$ for
$\zeta=\zeta_{3^m}$, so $h(\zeta)$ lies in the maximal real subfield
$K_m^+$ of $K_m=\mathbb Q(\zeta_{3^m})$, of index $2$. Hence
\begin{align*}
 \Res(\Phi_{3^m},h)
 &=\prod_{\zeta\ \mathrm{prim}}h(\zeta)
  =N_{K_m/\mathbb Q}\bigl(h(\zeta)\bigr)\\
 &=\Bigl(N_{K_m^+/\mathbb Q}\bigl(h(\zeta)\bigr)\Bigr)^{2}
  =P_m^2,\qquad P_m\in\ZZ .
\end{align*}
Since $x^{27}-1=\prod_{m=0}^{3}\Phi_{3^m}(x)$ and $\Phi_1(x)=x-1$,
equation~\eqref{eq:red} gives
$D=h(1)\,P_1^2P_2^2P_3^2$, which is \eqref{eq:val}.
Finally $3$ is totally ramified in $K_m$, with
$\pi_m=1-\zeta_{3^m}$ generating the prime above it and
$N_{K_m/\mathbb Q}(\pi_m)=\Phi_{3^m}(1)=3$; as
$h(\zeta)\equiv h(1)\pmod{\pi_m}$, we get $\pi_m\mid h(\zeta)$ exactly when
$3\mid h(1)$, and $v_3(P_m^2)=v_{\pi_m}(h(\zeta))$.
\end{proof}

\begin{proof}[Proof of Theorem \ref{thm:main}]
Both $h$ displayed in Theorem \ref{thm:main} follow from \eqref{eq:h} by
inspection, and both have $h(1)$ divisible by $3$; Lemma \ref{lem:sq} then
forces $v_3(P_m)\ge1$ for $m=1,2,3$. Since each $\Res(\Phi_{3^m},h)=9$ we have
$P_m=\pm3$ and $v_3(P_m)=1$, so \eqref{eq:val} reads
$v_3(D)=v_3(h(1))+6$, which is $1+6=7$ for W1 and $2+6=8$ for W2. The
displayed products $3\cdot9^3=2187$ and $9\cdot9^3=6561$ are the values
themselves, so nothing beyond $3$ divides them and the cofactor is $1$.
\end{proof}

\begin{remark}[a hand check of $P_1$]
Grouping the coefficients of $h$ by residue mod $3$ and using symmetry,
$h(\zeta_3)=A-B$ with $A=\sum_{i\equiv0}h_i$ and
$B=\sum_{i\equiv1}h_i=\sum_{i\equiv2}h_i$, so
$\Res(\Phi_3,h)=(A-B)^2$ and
\[
 P_1=h(1)-3H_1,\qquad
 H_1=\sum_{\substack{1\le i\le13\\ i\equiv\pm1\ (3)}}h_i .
\]
For W1, $H_1=h_1+h_4=2$ and $P_1=3-6=-3$; for W2,
$H_1=h_1+h_2+h_4+h_5=4$ and $P_1=9-12=-3$. Both give $P_1^2=9$.
\end{remark}

\begin{remark}
Exchanging $f$ and $g$ replaces $h$ by $-h$, and $M(-h)=-M(h)$ because $27$ is
odd; so $-3^7$ and $-3^8$ are measures as well.
\end{remark}

\section{Completeness}

\begin{theorem}\label{thm:complete}
Granting the necessity and achievability statements of Theorem 5.4 of
\cite{BP} at $k=3$, $p=3$, which are quoted and not reproved here,
\[
 \Sset(D_{54})=\bigl\{2^a3^bm : m\in\ZZ,\ \gcd(m,6)=1,\ a=0\text{ or }a\ge2,\
 b=0\text{ or }b\ge7\bigr\}.
\]
\end{theorem}

\begin{proof}
The inclusion $\subseteq$ is Theorem 5.4 of \cite{BP} at $k=3$: every measure
is of the form $2^a3^bm$ with $(m,6)=1$, $a=0$ or $a\ge2$, and $b=0$ or
$b\ge2\cdot3+1=7$.

For $\supseteq$, $D_G$ is multiplicative on $\ZZ[G]$, so $\Sset(G)$ is closed
under multiplication. From \cite{BP} (proof of Theorem 5.1) $\Sset(D_{54})$
contains every $2^am$ with $(m,6)=1$ and $a=0$ or $a\ge2$ --- these are the
$b=0$ members --- and every pure power $3^b$ with $b\ge3k=9$. Theorem
\ref{thm:main} adds the pure powers $3^7$ and $3^8$. The set of attainable
exponents $b$ therefore contains $\{0\}\cup\{7,8\}\cup\{b:b\ge9\}
=\{0\}\cup\{b:b\ge7\}$, and multiplying a pure power $3^b$ by a $b=0$ member
$2^am$ realises every $2^a3^bm$ of the stated form with $m>0$. Finally,
exchanging $f$ and $g$ replaces $h$ by $-h$ and hence $D$ by $-D$ (Remark at
the
end of Section \ref{sec:valuation}), so $\Sset(D_{54})$ is closed under
negation and the cases $m<0$ follow from the cases $m>0$.
\end{proof}

\begin{remark}
Measures with $v_3=7$ and $v_3=8$ but
nontrivial cofactor already follow from constructions printed in \cite{BP};
such cofactors cannot be removed afterwards, since $\Sset(G)$ is closed under
multiplication and not under division.
\end{remark}

\section{Reproduction}

Everything above is determined by the two pairs $(f,g)$ printed in Theorem
\ref{thm:main} and by \eqref{eq:red}. To redo the computation: form
$h=ff^*-gg^*$ in $\ZZ[x]/(x^{27}-1)$ and evaluate $D=\Res(x^{27}-1,h)$ ---
equivalently, the $27\times27$ circulant determinant of $h$, or the product
\[
 h(1)\,\Res(\Phi_3,h)\,\Res(\Phi_9,h)\,\Res(\Phi_{27},h),
\]
or the $54\times54$ determinant $\det(a_{uv^{-1}})$ built directly from the
multiplication table of $D_{54}$, which is the definition. All four routes agree
on both witnesses.

The accompanying program \texttt{verify.py} carries out the witness computations
in exact integer arithmetic: it builds $D_{54}$ from its presentation, computes
each determinant by all four routes, and confirms the valuations, the unit
cofactors and the factorisations $P_m^2=9$. It does not verify the necessity
statement or the infinite achievability results quoted from \cite{BP}.

\begin{thebibliography}{9}

\bibitem{BP}
M.~Boerkoel and C.~Pinner,
\emph{Minimal group determinants and the Lind--Lehmer problem for dihedral
groups},
Acta Arith. \textbf{186} (2018), no.~4, 377--395.
\href{https://doi.org/10.4064/aa180708-30-8}{doi:10.4064/aa180708-30-8};
e-print \href{https://arxiv.org/abs/1802.07336}{arXiv:1802.07336}.

\end{thebibliography}

\end{document}
